\section{Results}
\label{sec:results}
\subsection{Capability and formation}
The capability check returned 20 valid responses: all ten correct claims were
approved and all ten faulty claims rejected. These observations were excluded
from the study estimates. All five paired blocks completed, producing 100
formation verdicts and 240 recovery verdicts, with no missing or invalid scored
decisions. Independent replay and rescoring verified the saved evidence.

The reviewer rejected every formation claim under both peer policies. In the
prespecified comparison over episodes 2--10, approvals were 0/45 under each
policy, or 0/9 in every block. The first-episode control also had 0/5 approvals
per policy. Thus the study did not demonstrate the peer-induced unsupported
approval required to interpret later behavior as recovery from that failure.

\subsection{Outcomes under the recovery conditions}
Table~\ref{tab:recovery-observed} reports the full assigned schedule. In each
full-evidence condition, the reviewer rejected all 20 faulty claims and approved
all 20 correct claims. Every block had the same pattern: zero faulty approvals
and four correct approvals per condition. The primary B$-$A false-approval
difference was therefore zero in all five blocks.

\begin{table}[htbp]
\centering\small
\caption{Observed review outcomes. Counts combine five paired blocks, each with
four faulty and four correct claims per condition; the independent sample size
is five. All blocks had identical rates. Balanced accuracy (BA) averages correct
rejection and correct approval. Conditions are defined in Table~\ref{tab:arms}.}
\label{tab:recovery-observed}
\begin{tabularx}{\linewidth}{@{}lXrrr@{}}
\toprule
Arm & Condition & Faulty approved & Correct approved & BA (\%)\\
\midrule
A & Rejecting peer, full evidence & 0/20 & 20/20 & 100\\
B & Accepting peer, full evidence & 0/20 & 20/20 & 100\\
C & Rejecting peer, correction & 0/20 & 20/20 & 100\\
D & Accepting peer, correction & 0/20 & 20/20 & 100\\
E & Accepting peer, correction and reset & 0/20 & 20/20 & 100\\
F & Accepting peer, restricted evidence & 0/20 & 0/20 & 50\\
\bottomrule
\end{tabularx}
\end{table}

Under continued evidence restriction (F), the reviewer rejected all 40 claims.
Compared with F, B had an observed correct-approval difference of 100 percentage
points and a balanced-accuracy difference of 50 points; false approval remained
zero. These are descriptive differences in useful task performance after the
same assigned history, not evidence of repairing an induced harmful convention.
Rejection under F was consistent with the instruction to reject without records.

Correction, added history removal, and the correction-by-history interaction
all had observed differences of zero in false approval, correct approval, and
balanced accuracy. Full-evidence performance was already at ceiling without
these interventions, leaving no observed errors for them to repair.

\subsection{Uncertainty and interpretation}
All paired effects were constant across the five blocks. The prespecified
analysis therefore returned no informative bootstrap, Student-$t$, or combined
intervals. Because there were no missing decisions, the missing-outcome bounds
equaled the observed estimates. Those bounds describe unresolved outcomes in
this sample; they do not bound population uncertainty. We make no equivalence,
general absence-of-effect, or intervention-efficacy claim.
